\documentclass[a4paper]{article}
% Español (México) con XeLaTeX: fontspec para fuentes Unicode, polyglossia
% para la división silábica y los nombres del documento en la variante mexicana.
\usepackage{fontspec}
\usepackage{polyglossia}
\setdefaultlanguage[variant=mexican]{spanish}
% Los nombres chinos van como «pinyin (original)»: xeCJK da a los caracteres
% Han su propia fuente; sin ella XeLaTeX los omite con un aviso.
% FandolSong no cubre algunos caracteres de nombres (U+73FA); WenQuanYi Zen Hei sí.
\usepackage[AutoFallBack=true]{xeCJK}
\setCJKmainfont{FandolSong-Regular.otf}
\setCJKfallbackfamilyfont{\CJKrmdefault}{WenQuanYi Zen Hei}
% polyglossia no traduce los nombres de listings.
\AtBeginDocument{\providecommand{\lstlistingname}{}\renewcommand{\lstlistingname}{Listado}%
  \providecommand{\lstlistlistingname}{}\renewcommand{\lstlistlistingname}{Índice de listados}}
\usepackage{amsmath, amssymb}
\usepackage{graphicx}
\usepackage[margin=2.5cm]{geometry}
\usepackage[most]{tcolorbox}
\usepackage{etoolbox}
\usepackage{listings}
\usepackage{hyperref}
\usepackage{booktabs}
\usepackage{subcaption}
\usepackage{float}

\newtcolorbox{knowledgebox}[1]{enhanced, colback=blue!5!white, colframe=blue!75!black, colbacktitle=blue!75!black, coltitle=white, fonttitle=\bfseries, title={#1}, attach boxed title to top left={yshift=-2mm, xshift=2mm}, boxrule=1pt, sharp corners}
\newtcolorbox{importantbox}[1]{enhanced, colback=yellow!10!white, colframe=yellow!80!black, colbacktitle=yellow!80!black, coltitle=black, fonttitle=\bfseries, title={#1}, sharp corners}
\newtcolorbox{warningbox}[1]{enhanced, colback=red!5!white, colframe=red!75!black, colbacktitle=red!75!black, coltitle=white, fonttitle=\bfseries, title={#1}, sharp corners}

\lstset{language=Python, basicstyle=\ttfamily\small, keywordstyle=\color{blue}, stringstyle=\color{red!60!black}, commentstyle=\color{green!60!black}, breaklines=true, frame=single, numbers=left, numberstyle=\tiny\color{gray}, captionpos=b, extendedchars=true}

\newcommand{\notetitle}{Gemini 2.5 Pro + Agent = medalla de oro en la IMO: Actor-Critic Workflow}
\newcommand{\noteauthors}{Elaboradas a partir de materiales públicos del curso}
\newcommand{\notedate}{2025}
\newcommand{\videochannel}{Wudaokou Nashi (五道口纳什)}
\newcommand{\videopublishdate}{2025}
\newcommand{\videoduration}{aproximadamente 25 minutos}
\newcommand{\videourl}{https://www.bilibili.com/video/BV1Aq8bz1Emj}
\newcommand{\videocoverpath}{cover.jpg}
\newcommand{\repourl}{https://github.com/hqhq1025/ai-course-notes}


% Footer with repo info
\usepackage{fancyhdr}
\pagestyle{fancy}
\fancyhf{}
\fancyfoot[C]{\thepage}
\fancyfoot[R]{\footnotesize\color{gray}\href{\repourl}{AI Course Notes}}
\renewcommand{\headrulewidth}{0pt}
\renewcommand{\footrulewidth}{0pt}

\begin{document}
\begin{titlepage}
\centering
{\Large Notas del curso\par}\vspace{1.2cm}
{\huge\bfseries \notetitle\par}\vspace{0.8cm}
{\large \noteauthors\par}\vspace{0.3cm}
{\large \notedate\par}\vspace{1.2cm}
\ifdefempty{\videocoverpath}{}{\includegraphics[width=0.82\textwidth,height=0.45\textheight,keepaspectratio]{\videocoverpath}\par}
\vfill
\begin{tcolorbox}[width=0.9\textwidth, colback=black!2!white, colframe=black!60, sharp corners]
\textbf{Autor o canal del video}: \videochannel\par\textbf{Fecha de publicación}: \videopublishdate\par\textbf{Duración del video}: \videoduration\par
\ifdefempty{\videourl}{}{\textbf{Enlace del video}: \href{\videourl}{\nolinkurl{\videourl}}\par}\par
\textbf{Página del proyecto}: \href{\repourl}{\nolinkurl{\repourl}}\par
\end{tcolorbox}
\end{titlepage}
\tableofcontents\newpage

\section{Contexto del trabajo y hallazgo central}

\subsection{Un resultado alentador}

En esta sesión se presenta un trabajo que resultó sumamente inspirador para el autor: \textbf{Gemini 2.5 Pro} nativo (la versión sin el tune oficial de DeepMind), mediante un Agent Workflow cuidadosamente diseñado, resolvió los primeros cinco de los seis problemas de la IMO 2025, alcanzando el nivel de medalla de oro humano.

\begin{importantbox}{Hallazgo central}
La palabra clave es \textbf{capable}: el modelo en sí tiene la capacidad de resolver estos problemas matemáticos extremadamente difíciles; el desempeño deficiente anterior se debía a que no se había usado un Workflow adecuado para maximizar el potencial del modelo. Esto tiene implicaciones enormes para el diseño de agentes.
\end{importantbox}

\subsection{Resumen de la sección}

Este trabajo demuestra la importancia del diseño del Agent Workflow: con el mismo modelo, un Workflow distinto puede traer un salto cualitativo.

\section{Proceso completo del workflow}

\subsection{Proceso de seis pasos}

\begin{enumerate}
    \item \textbf{Paso 1: solución inicial}——produce una solución inicial a partir del prompt y de Gemini 2.5 Pro
    \item \textbf{Paso 2: Self-Improvement}——comparte el contexto del paso 1 y continúa haciendo refine (porque el paso 1 puede terminar de manera apresurada por agotarse el Thinking Budget)
    \item \textbf{Paso 3: Verification}——verifica la solución del paso 2 y produce un verdict (válido/inválido) y un bug report
    \item \textbf{Paso 4: Bug Report}——extrae la ubicación y la causa concretas del error
    \item \textbf{Paso 5: Correction}——corrige la solution con base en el bug report
    \item \textbf{Paso 6: ciclo}——los pasos 3 a 5 se repiten en ciclo
\end{enumerate}

\begin{importantbox}{Condiciones de terminación}
\begin{itemize}
    \item \textbf{Aceptación}: una solution pasa la verification \textbf{5 veces} consecutivas $\rightarrow$ se considera resuelto el problema
    \item \textbf{Rechazo}: tras \textbf{10} ciclos sigue sin pasar $\rightarrow$ se abandona la dirección actual y se vuelve a hacer sample
\end{itemize}
\end{importantbox}

\subsection{Patrón Evaluator-Optimizer}

\begin{knowledgebox}{Correspondencia con el patrón clásico de Building Effective Agents}
Este workflow es, en esencia, el patrón \textbf{Evaluator-Optimizer}:
\begin{itemize}
    \item el Generator (generador) produce la solution
    \item el Evaluator (evaluador) la verifica y da feedback
    \item si se acepta, se produce como salida; si se rechaza, se vuelve a generar con base en el feedback
\end{itemize}
\end{knowledgebox}

\subsection{El papel decisivo de la calidad de la solución inicial}

\begin{warningbox}{La clave del éxito está en la solución inicial}
Si la solución inicial que se genera en el paso 2 tiene un overlap considerable con la solución correcta, entonces el ciclo posterior de verificación y corrección tiene una alta probabilidad de resolver el problema. Pero si la dirección de la solución inicial ya está equivocada (se desvía demasiado), es poco probable que las correcciones posteriores puedan salvarla.
\end{warningbox}

\subsection{Resumen de la sección}

Todo el workflow es un ciclo iterativo estándar de generación, verificación y corrección, cuyas condiciones de terminación se basan en el conteo de veces consecutivas que se pasa o no se pasa la verificación.
\section{Análisis del Prompt del Paso 1: generación de la solución inicial}

\subsection{Estructura del Prompt}

El System Prompt del Paso 1 es un texto en Markdown estructurado estándar, dividido en tres partes:

\begin{enumerate}
    \item \textbf{Instrucciones centrales}:
    \begin{itemize}
        \item El rigor es fundamental (la mayoría de los problemas de la IMO son de demostración, y el criterio final es el rigor lógico)
        \item Sé honesto respecto al resultado: no lo fuerces (para evitar el Reward Hacking)
        \item Usa LaTeX en todo el texto
    \end{itemize}
    \item \textbf{Formato de salida (Adaptive Prompt)}:
    \begin{itemize}
        \item Summary: incluye el Verdict (si se encontró una solución completa) y un resumen del método
        \item Detailed Solution: la solución detallada
    \end{itemize}
    \item \textbf{Instrucciones de Self-Verification}
\end{enumerate}

\subsection{Intención de diseño del formato de salida}

\begin{importantbox}{El Verdict se usa para decidir la ramificación}
El Verdict dentro del Summary no solo presenta información, sino que además es \textbf{el criterio para decidir la ramificación} del Workflow posterior: si el Verdict es false (no se encontró una solución completa), se omite la verificación posterior y se vuelve a hacer sample directamente. La Detailed Solution, en cambio, se envía a la etapa de Verification.
\end{importantbox}

El código de este texto usa \textbf{expresiones regulares} para extraer estas salidas estructuradas (en lugar de una API de salida estructurada), porque la capacidad de seguimiento de instrucciones de Gemini 2.5 Pro es lo bastante sólida.

\subsection{Paso 2: Self-Improvement}

El Paso 2 \textbf{comparte el contexto} con el Paso 1. Su Prompt es sencillo: «tienes una oportunidad de mejorar tu resultado: haz un review cuidadoso de tu solution, corrige errores y completa los justification gap».

La razón para introducir el Paso 2 es que, en el Paso 1, el Thinking Budget (32768 tokens) puede agotarse por completo, lo que obliga al modelo a terminar de pensar de forma apresurada. El Paso 2 le da al modelo la oportunidad de seguir pensando.

\subsection{Resumen de la sección}

Lo esencial del diseño del Prompt es que la salida estructurada sienta las bases para decidir la ramificación y el flujo de datos del Workflow posterior.
\section{Análisis del Verification Prompt}

\subsection{El papel y las instrucciones del Verifier}

El System Prompt de Verification también es Markdown estructurado y contiene:

\begin{enumerate}
    \item \textbf{Interpretación de un papel}: eres un matemático excepcional, un revisor con el rigor de nivel IMO
    \item \textbf{Indicación central}: \textbf{eres solo el Verifier, no el Solver}: no intentes corregir errores ni llenar gap, solo verifica step by step
    \item \textbf{Clasificación de problemas}:
    \begin{itemize}
        \item \textbf{Critical Error}: error lógico grave (por ejemplo, de $A > B, C > D$ deducir $A-C > B-D$)
        \item \textbf{Gap}: paso omitido, ruptura lógica
    \end{itemize}
    \item \textbf{Formato de salida}: Summary (con Verdict: valid/invalid) + List of Findings (bug report) + Detailed Verification Log
\end{enumerate}

\begin{warningbox}{El Verifier no debe intentar resolver el problema}
El Verification Prompt exige explícitamente que el Verifier no intente corregir errores ni llenar gap, y que solo verifique. Esta \textbf{separación de papeles} es clave en el diseño del Workflow: el Generator y el Verifier cumplen cada uno su función.
\end{warningbox}

\subsection{Resumen de la sección}

Mediante una definición estricta de papeles y un formato de salida estructurado, el Verification Prompt garantiza que el Verifier solo verifique y no resuelva; su Verdict y su Bug Report resultantes impulsan la Correction posterior y la decisión de terminar el ciclo.

\section{Correction y el ciclo iterativo}

\subsection{El paso de Correction}

Cuando el Verifier determina que la solution no es válida, se extrae el bug report y se envía al paso de Correction:
\begin{itemize}
    \item Entrada: el problema original + la solution actual + el bug report
    \item Prompt: ``Below is a bug report. If you agree with it, fix the issues. Your output should follow the system prompt instructions.''
    \item Salida: la nueva solution corregida
\end{itemize}

\subsection{El proceso iterativo completo}

Tomando el primer problema como ejemplo:
\begin{enumerate}
    \item Exploración inicial: Step 1 + Step 2 generan la solución inicial
    \item Verificación: el Verifier determina No (hay un bug)
    \item Corrección: se genera una nueva solución a partir del bug report $\rightarrow$ verificación No $\rightarrow$ corrección $\rightarrow$ verificación No
    \item Tras la cuarta corrección se pasa la verificación: correct count = 1
    \item Se pasa la verificación 5 veces consecutivas y finalmente se acepta
    \item En total se pasa por 8 iteraciones, y el ciclo externo solo tiene 1 run
\end{enumerate}

\subsection{Resumen de la sección}

Todo el ciclo iterativo muestra la fuerza del patrón Actor-Critic: aunque la solución inicial tenga varios errores, mediante el ciclo repetido de verificación y corrección finalmente se puede converger a la solución correcta.

\section{Aspectos a tener en cuenta al usar la API de Gemini}

\begin{knowledgebox}{Detalles técnicos}
\begin{itemize}
    \item Los roles de la API de Gemini son system / user / model (no system / user / assistant)
    \item La salida de Thinking es el resultado después de un summarize, no el proceso de thinking nativo
    \item El Thinking Budget se configura en 32768 tokens, pero a veces llega a más de 34000 en la práctica (el control no es estricto)
    \item Como el proceso de Thinking es muy largo (decenas de miles de tokens), usar HTTP Request produce fácilmente un timeout; se recomienda usar el SDK con Streaming
\end{itemize}
\end{knowledgebox}
\section{Síntesis y ampliación}

\begin{enumerate}
    \item El \textbf{ciclo Generate-Verify-Correct} es un patrón de Workflow poderoso y general; el éxito de este artículo en los problemas de la IMO valida su eficacia.
    \item \textbf{La calidad de la solución inicial determina el resultado}: si la dirección de la solución inicial es correcta (con un overlap considerable respecto a la solución estándar), el ciclo posterior probablemente pueda corregir los detalles; si la dirección es incorrecta, es difícil corregirla.
    \item \textbf{El diseño del Prompt es el alma del Workflow}: un formato de salida estructurado sienta la base para la ramificación automatizada de decisiones y el flujo de datos.
    \item \textbf{Separación de roles}: el Generator solo se encarga de generar, el Verifier solo de verificar y el Corrector solo de corregir; que cada quien cumpla su función es el principio central del diseño Multi-Agent.
    \item \textbf{El valor del diseño de Agent}: el mismo modelo (Gemini 2.5 Pro), mediante el diseño del Workflow, puede pasar de «no poder resolverlo» a «nivel de medalla de oro».
\end{enumerate}

\end{document}
